\documentclass[11pt,a4paper]{article}
\usepackage[utf8]{inputenc}
\usepackage[english,ukrainian]{babel}
\usepackage{amsmath,amssymb,amsfonts}
\usepackage{hyperref}
\usepackage{geometry}
\geometry{margin=1in}

\title{\textbf{Семантична інженерія архітектури інформації:\\ фрактальний простір об'єктивних сенсів та когнітивний екзоскелет}}
\author{Ваше Ім'я / Незалежний дослідник}
\date{\today}

\begin{document}

\maketitle

\begin{abstract}
Сучасні інформаційні системи створюють надмірне когнітивне навантаження та інформаційний хаос через пасивне накопичення даних у бібліотеках. У цій роботі представлено новий концептуальний фреймворк семантичної інженерії інформації (СІІ). На відміну від статичних архітектур, пропонований підхід стверджує, що інформація має не накопичуватися, а через резонанс і сигнал сприйняття активувати дію. Впровадження фрактальної мапи ранжування та «кристаликів сенсу» усуває застарілу концепцію «навчання», перетворюючи читача на фахівця через інтуїтивну GPS-навігацію. Архітектурну модель продемонстровано на прикладі контуру адміністративно-фінансових сенсів (КЕП-Дія/ФОП — Банківський рахунок).
\end{abstract}

\section{Анонс (Таймлайн статті / Roadmap)}
У цій роботі розкривається перехід від інформаційних бібліотек до систем прямої дії. Стаття побудована за таким таймлайном:
\begin{enumerate}
    \item \textbf{Концепт-контент (Парадигма активації):} Чому інформація дорівнює дії, а не накопиченню фактів, і як працює «кристалик сенсу».
    \item \textbf{Архітектурний приклад:} Практичний контур взаємодії (КЕП-Дія/ФОП — Банківський рахунок).
    \item \textbf{Динамічний глосарій:} Базові терміни семантичної інженерії та їхнє розширення.
    \item \textbf{Методологічні межі:} Чому глобальна архітектура завжди незавершена, а локальна сенсова візуалізація завершена по темі.
    \item \textbf{Кінцева заставка (Вихід у дію):} Перетворення читача на фахівця без етапу заучування.
\end{enumerate}

\section{Концепт-Контент: Інформація = Дія}
Експоненційне зростання цифрових даних перенесло вузьке місце продуктивності на ємність когнітивної обробки. Традиційні парадигми вимагають виснажливого запам'ятовування. Наша концепція доводить: інформація не повинна накопичуватися у бібліотеці. Її мета — через резонанс і сигнал сприйняття активувати дію. 

Фундаментальним елементом є виключення хаосу через фрактальну організаційну структуру — від макрорівня архітектури до мікрорівня «кристалика сенсу». Наприклад, КЕП не є сухим фактом; це заряджений вузол: \textit{особиста відповідальність і право, а також можливості і вимоги}. Це фрактальне зерно миттєво запускає дію без заучування.

\section{Архітектурний приклад: Контур «КЕП-Дія/ФОП — Банківський рахунок»}
Замість вивчення інструкцій із цифрової бібліотеки, інформаційний простір створює сигнал сприйняття. Статус ФОП (КЕП-Дія) і банківський рахунок об'єднуються у прозору структуру, яка розгортається в розумовому просторі, миттєво переводячи користувача в режим прямої кваліфікованої дії (нульова фаза взаємодії).

\section{Динамічний глосарій семантичної інженерії}
\begin{itemize}
    \item \textbf{Семантична інженерія інформації (СІІ):} Пряме проєктування архітектури знань, що замінює пасивне накопичення бібліотек на активну генерацію сенсів.
    \item \textbf{GPS-навігація по темі:} Архітектурний принцип взаємодії без потреби заучування фактів.
    \item \textbf{Резонанс і сигнал сприйняття:} Механізм активації, де архітектура даних миттєво ініціює цілеспрямовану дію.
    \item \textbf{Когнітивний екзоскелет:} Сукупність зовнішніх структур, що збільшують швидкість переходу від сприйняття до дії.
\end{itemize}

\section{Кінцева заставка: Відкрита методологія та еволюція глосарію}
Будь-яка архітектура семантичної інженерії є принципово відкритою та незавершеною на рівні глобальної методології. Проте завдяки фрактальній природі об'єктивна сенсова візуалізація досягає локальної завершеності по конкретній темі. Базова структура постійно розширюється через нові блоки та сенсовий глосарій. Незалежно від обраної форми передачі (слова, терміни, карти, відео, таблиці), інформаційний об'єкт діє лише як тригер, що інтерпретується особистістю через її унікальне бачення.

\end{document}
